\documentclass[12pt,a4paper]{article}

\usepackage{amssymb}
\usepackage{amsmath}
\usepackage{amsthm}
\usepackage{bm}
\usepackage{enumitem}
\usepackage{mathtools}
\usepackage{lipsum}
\usepackage[brazil]{babel}
\usepackage[utf8]{inputenc}
\usepackage[T1]{fontenc}
\usepackage[pdftex]{graphicx}
\usepackage{color}
\usepackage{tikz}
\usetikzlibrary{arrows, calc, external}
\usepackage{verbatim}
\usepackage{url}
\usepackage{longtable}
\usepackage{geometry}
\geometry{
    a4paper,
    margin=2cm
}

\usepackage{helvet}
\renewcommand{\familydefault}{\sfdefault}

\usepackage{hyperref} 
\usepackage[alf]{abntex2cite}
\providecommand{\abntnextkey}{}

\usepackage{titlesec}
\usepackage{setspace}

\newtheorem{prop}{Proposição}

\begin{document}
\onehalfspacing

\begin{center}
    \textbf{Fundação Universidade Federal do ABC}\\
    \textbf{Pró-reitoria de pesquisa}\\
    \small Av. dos Estados, 5001, Santa Terezinha, Santo André/SP, CEP 09210-580\\
    Bloco L, 3º Andar, Fone (11) 3356-7617 | iniciacao@ufabc.edu.br
\end{center}

\vspace{0.8cm}

\noindent \textbf{Relatório Parcial de Iniciação Científica referente ao Edital:} 01/2025 PIC\slash PIBIC\slash PIBITI\slash PIBIC-AF \\[0.5cm]
\noindent \emph{Na capa deverá conter tanto a assinatura do aluno quanto a assinatura do orientador}\\[0.5cm]
\noindent \textbf{Nome do aluno:} \hrulefill \\
\noindent \textbf{Assinatura do aluno:} \hrulefill \\[0.5cm]
\noindent \textbf{Nome do orientador:} \hrulefill \\
\noindent \textbf{Assinatura do orientador:} \hrulefill \\[0.5cm]
\noindent \textbf{Título do projeto:} Problemas de Fluxos em Redes: Teoria, Algoritmos e Implementações \\[0.5cm]
\noindent \textbf{Palavras-chave do projeto:} Otimização Combinatória, Teoria dos Grafos, Fluxos \\[0.5cm]
\noindent \textbf{Área do conhecimento do projeto:} Ciência da Computação, Matemática Discreta \\[0.5cm]
\noindent \textbf{Bolsista:} ( \textbf{\textsf{X}} ) Sim \quad ( \quad ) Não. \\
\textbf{Em caso positivo, informar a modalidade da bolsa:} PIBIC - CNPq \\[1cm]

\begin{center}
    \textbf{Santo André}\\
    \textbf{\today}
\end{center}

\newpage
\renewcommand{\contentsname}{Sumário}
\tableofcontents
\newpage

\section{Resumo}
Fluxos em digrafos constituem uma área fundamental da Otimização Combinatória, com amplas aplicações práticas em redes de transporte e alocação de recursos. Este projeto adota uma abordagem que combina fundamentos teóricos, implementação algorítmica e validação experimental para o estudo de problemas de fluxo máximo e fluxo de custo mínimo~\cite{ahuja1993network}. Nesta etapa parcial, foram estudados algoritmos de fluxo máximo e técnicas de modelagem. A validação prática está sendo realizada através de experimentos computacionais utilizando instâncias de \textit{benchmarks} reconhecidas~\cite{Jensen2022}.

\section{Introdução}
A Otimização Combinatória é uma área de pesquisa situada na interseção entre Matemática e Ciência da Computação~\cite{korte2018combinatorial}. Do ponto de vista matemático, a área mantém forte influência com a Teoria dos Grafos, uma vez que grafos constituem objetos combinatórios com amplo poder de modelagem~\cite{bondy2008graph, diestel2017graph}. 
Sob a perspectiva computacional, o foco reside na obtenção de algoritmos eficientes que produzam soluções garantidamente ótimas~\cite{schrijver2003combinatorial, cook1998combinatorial}. Dentre os tópicos de destaque, os problemas de fluxos em digrafos ocupam posição central, tratando de redes onde recursos devem ser transportados respeitando restrições de capacidade~\cite{ahuja1993network}.
O objetivo deste projeto é proporcionar uma formação sólida na área, integrando teoria e prática. Até o presente momento, as atividades focaram na consolidação teórica e na implementação de algoritmos para o problema do fluxo máximo.

\section{Fundamentação teórica}
O problema do $st$-fluxo máximo, onde se busca maximizar a quantidade de fluxo enviada de um vértice fonte $s$ para um vértice de destino $t$, é um dos pilares deste estudo~\cite{ford1956maximal}. Foram analisados algoritmos como os de Ford-Fulkerson, Edmonds-Karp e Dinic~\cite{ford1956maximal, edmonds1972theoretical, dinic1970algorithm}. 
O projeto também contempla o problema de fluxo de custo mínimo, que busca minimizar o custo total enquanto satisfaz as demandas da rede. O algoritmo principal a ser implementado para este fim é o \textit{network simplex}, uma especialização do algoritmo simplex~\cite{dantzig1951application, ahuja1993network}. 
A fundamentação teórica revisada nesta etapa reafirma que muitos problemas combinatórios, como o emparelhamento máximo em grafos bipartidos, podem ser resolvidos via reduções para problemas de fluxo~\cite{ahuja1993network}.

\section{Metodologia}
\subsection{Materiais e Métodos}
A metodologia baseia-se em ciclos de estudo bibliográfico seguidos de implementação. Os algoritmos estão sendo desenvolvidos em linguagem C++ para garantir alta performance, utilizando listas de adjacência otimizadas. 
Os experimentos serão conduzidos em instâncias padronizadas das coleções DIMACS~\cite{Jensen2022}, permitindo a comparação de desempenho. A análise baseia-se em métricas de tempo de execução e corretude dos algoritmos implementados frente às soluções ótimas conhecidas~\cite{ford1956maximal}.

\subsection{Etapas da pesquisa}
As atividades do projeto foram estruturadas e encontram-se no seguinte status:
\begin{itemize}
    \item \textbf{Atividade 1:} Estudo dos fundamentos de Otimização Combinatória e Teoria dos Grafos~\cite{cook1998combinatorial}. \textbf{[Concluída]}
    \item \textbf{Atividade 2:} Estudo teórico do problema de fluxo máximo e algoritmos clássicos~\cite{ford1956maximal, edmonds1972theoretical, dinic1970algorithm, goldberg1988new}. \textbf{[Concluída]}
    \item \textbf{Atividade 3:} Implementação em C++ dos algoritmos de fluxo máximo. \textbf{[Concluída]}
    \item \textbf{Atividade 4:} Estudo de técnicas de modelagem e redução de problemas. \textbf{[Concluída]}
    \item \textbf{Atividade 5:} Implementação de problemas modelados (Emparelhamento e Caminhos Disjuntos). \textbf{[Em andamento]}
    \item \textbf{Atividade 6:} Elaboração do relatório parcial (este documento). \textbf{[Em andamento]}
    \item \textbf{Atividade 7:} Estudo teórico do problema de fluxo de custo mínimo~\cite{cook1998combinatorial}. \textbf{[A iniciar]}
    \item \textbf{Atividade 8:} Implementação em C++ do algoritmo \textit{network simplex}. \textbf{[A iniciar]}
    \item \textbf{Atividade 9:} Testes experimentais em instâncias de larga escala. \textbf{[A iniciar]}
    \item \textbf{Atividade 10:} Análise comparativa de desempenho dos algoritmos. \textbf{[A iniciar]}
    \item \textbf{Atividade 11:} Elaboração do relatório final e artigo de conclusão. \textbf{[A iniciar]}
\end{itemize}

\section{Resultados parciais}
Nesta seção, devem ser apresentados os tempos de execução obtidos pelos algoritmos de Edmonds-Karp e Dinic nas primeiras instâncias testadas, bem como a verificação da corretude através do teorema do corte mínimo~\cite{ford1956maximal}.

\section{Conclusões preliminares e próximos passos}
Os objetivos propostos para a primeira metade do cronograma foram alcançados, com a implementação bem-sucedida dos algoritmos de fluxo máximo. Os próximos passos focam na transição para problemas de fluxo de custo mínimo~\cite{ahuja1993network} e na execução da bateria completa de testes experimentais.

\newpage
\bibliographystyle{abntex2-alf}
\bibliography{cit}

\end{document}
